\documentclass[11pt,a4paper]{amsart}

\usepackage[margin=1in]{geometry}
\usepackage[T1]{fontenc}
\usepackage{lmodern}
\usepackage{microtype}
\usepackage{amsmath,amssymb,amsthm,mathtools}
\usepackage{booktabs,tabularx}
\usepackage[colorlinks=true,linkcolor=blue,urlcolor=blue]{hyperref}

\newtheorem{theorem}{Theorem}[section]
\newtheorem{proposition}[theorem]{Proposition}
\newtheorem{corollary}[theorem]{Corollary}
\theoremstyle{remark}
\newtheorem{remark}[theorem]{Remark}

\newcommand{\F}{\mathbb F}
\newcommand{\Z}{\mathbb Z}
\newcommand{\wt}{\operatorname{wt}}
\newcommand{\ord}{\operatorname{ord}}
\newcommand{\Span}{\operatorname{span}}
\newcommand{\im}{\operatorname{im}}

\title{An Admissible-Branch Correction for the Rate-Half F2 Kernel}
\author{Experimental theorem packet}
\date{August 2026}

\begin{document}
\begin{abstract}
We audit a rate-half proof route at the maximal smooth order $n=2^{41}$.
The frequently used formula
\(\ord_n(p)=2^{(41-v_2(p-1))_+}\) applies to
$p\equiv1\pmod4$, not to every admissible odd characteristic.  The official
row $p=2^{61}-1$, $q=p^2$ is a generating counterexample to the resulting
bounded proportionality-class census.  We give the corrected
$p\equiv3\pmod4$ order law and replace the false direct-sum model by one
Frobenius-coupled negacyclic root code with exact rank.  We also record the
pointwise weighted-collision floor and the correct Fourier factor
$1+\cos$, which differs from the unweighted ternary factor $1+2\cos$.
No MCA numerator, adjacent row, or Prize endpoint is closed here.
\end{abstract}
\maketitle

\section{Scope and notation}

Fix
\[
 n=2^{41},\qquad q=p^e<2^{256},\qquad n\mid q-1,
\]
where $p$ is odd.  Write $k=\ord_n(p)$.  The field cap and the usual
lifting-the-exponent calculation give
\[
 e\le 6,\qquad k\mid e,\qquad k\in\{1,2,4\}.
\]
We consider the first $R$ odd moment equations on an antipodal half-window.
The deployed range satisfies
\[
 2R<2^{36};
\]
for example this follows from $2R\le \lceil n/41\rceil+1$.

For a linear moment map $A:\F_p^m\to V$, put
\[
 K=\ker A,
 \qquad
 Z(A)=\sum_{\epsilon\in K\cap\{-1,0,1\}^m}2^{-\wt(\epsilon)}.
\]

\section{The omitted residue class}

\begin{proposition}[Correct dyadic order laws]\label{prop:order}
For $a\ge2$ and odd $p$:
\[
 \ord_{2^a}(p)=
 \begin{cases}
  2^{(a-v_2(p-1))_+},&p\equiv1\pmod4,\\[2mm]
  2^{\max\{1,a-v_2(p+1)\}},&p\equiv3\pmod4.
 \end{cases}
\]
Consequently, on an admissible minus-branch row
$p\equiv3\pmod4$, if $b=v_2(p+1)$, then
\[
 b\ge39,
 \qquad
 k=2^{\max\{1,41-b\}}\in\{2,4\}.
\]
\end{proposition}

\begin{proof}
For $p\equiv1\pmod4$, repeated squaring and LTE give
$v_2(p^{2^j}-1)=v_2(p-1)+j$.  For $p\equiv3\pmod4$, the first square gives
$v_2(p^2-1)=1+v_2(p+1)$ and each further square raises the valuation by one.
The least power of two whose exponent reaches $a$ gives the two formulas.
Since $k$ is a power of two dividing $e\le6$, it belongs to
$\{1,2,4\}$; the minus branch has $k\ge2$, yielding the final claim.
\end{proof}

\begin{theorem}[Official generating counterexample]\label{thm:m61}
Let
\[
 p=2^{61}-1=2305843009213693951,
 \qquad q=p^2.
\]
Then $p$ is prime, $q<2^{256}$, and
\[
 p\equiv-1\pmod{2^{41}},
 \qquad
 \ord_{2^{41}}(p)=2=e.
\]
Thus this is an admissible generating row.  Nevertheless
\[
 \left|\mu_{2^{41}}\cap\F_p^*\right|
 =\gcd(2^{41},p-1)=2.
\]
After one representative is retained from every antipodal pair, all
$2^{40}$ positions are distinct $\F_p$-proportionality classes.  In
particular, the all-admissible claim that there are at most four such classes
is false.
\end{theorem}

\begin{proof}
The Lucas--Lehmer test gives a primality certificate for the Mersenne number
$2^{61}-1$; the companion verifier replays it.  The congruence
$p\equiv-1\pmod{2^{41}}$ proves both $2^{41}\mid p^2-1$ and order two.
Also $p^2<2^{122}<2^{256}$.  Two roots are proportional over $\F_p$ exactly
when their quotient lies in the displayed intersection.  Its only elements
are $\{1,-1\}$, and the antipodal quotient has already removed the second
possibility.
\end{proof}

\begin{remark}[Surviving plus branch]
The counterexample does not affect the bounded-class reduction on
$p\equiv1\pmod4$.  In that branch the first formula in
Proposition~\ref{prop:order}, together with $k\mid e\le6$, gives
$k\in\{1,2,4\}$ and the established prime-field class decomposition.  The
repair is a branch split, not a rejection of the plus-branch theorem.
\end{remark}

\section{The coupled minus-branch code}

Let $W$ be either the exact-order or nested window at order $2^a$, with
$a\in\{40,41\}$, and let $m=|W|/2$.  There is an element $\omega$ of order
$2m$ for which a coefficient vector
$\epsilon=(\epsilon_0,\ldots,\epsilon_{m-1})$ has coefficient polynomial
\[
 P_\epsilon(X)=\sum_{s=0}^{m-1}\epsilon_sX^s
\]
and the first $R$ odd moments vanish exactly when
\[
 P_\epsilon(\omega^{2j-1})=0,
 \qquad 1\le j\le R.                                  \tag{3.1}
\]
For a nested window take $\omega$ to be a generator of $W$.  For an
exact-order window, write its representatives as
$y^{2s+1}$ and take $\omega=y^2$; the discarded factor $y^{2j-1}$ is
nonzero.

\begin{theorem}[Frobenius-closed negacyclic reduction]\label{thm:coupled}
Assume $p\equiv3\pmod4$ and the admissible hypotheses above.  Put
\[
 h=
 \begin{cases}
  2,&W\text{ exact-order at }2^{40}\text{ or }2^{41},\\
  2,&W\text{ nested at }2^{40},\\
  k,&W\text{ nested at }2^{41}.
 \end{cases}
\]
Then the Frobenius-closed root set
\[
 \Omega_W=
 \left\{\omega^{p^i(2j-1)}:
 0\le i<h,\ 1\le j\le R\right\}
\]
has exactly $hR$ elements.  Its root polynomial
\[
 G_W(X)=\prod_{\alpha\in\Omega_W}(X-\alpha)
\]
lies in $\F_p[X]$, and the moment kernel is exactly
\[
 K_W=
 \left\{\operatorname{coeff}(P):
 P\in\F_p[X],\ \deg P<m,\ G_W\mid P\right\}.          \tag{3.2}
\]
Consequently
\[
 \operatorname{rank}_{\F_p}A_W=hR,
 \qquad
 \dim_{\F_p}K_W=m-hR.                                 \tag{3.3}
\]
\end{theorem}

\begin{proof}
The relevant Frobenius order is $\ord_{2m}(p)$, which is the printed $h$ by
Proposition~\ref{prop:order}.  The smallest frequency modulus is
$2m=2^{39}$, whereas all selected exponents are less than
$2R<2^{36}$.  If $v_2(p+1)$ reaches the modulus exponent, Frobenius sends a
small odd exponent to its negative.  At the next modulus it sends it to a
half-period translate of its negative.  In the unique order-four case the
four conjugates occupy the four quarter-period bands.  These bands are
disjoint because their widths are below one eighth of the smallest modulus.
Hence $|\Omega_W|=hR$.

The set is Frobenius-stable, so $G_W\in\F_p[X]$.  A polynomial over $\F_p$
vanishing at the roots in (3.1) vanishes at all their Frobenius conjugates,
and the converse is immediate.  Distinct-root divisibility gives (3.2).
Multiplication by the monic degree-$hR$ polynomial $G_W$ identifies
polynomials of degree less than $m-hR$ with the kernel, proving (3.3).
\end{proof}

\begin{corollary}[Signed distance]\label{cor:distance}
Every nonzero
$\epsilon\in K_W\cap\{-1,0,1\}^m$ has
\[
 \wt(\epsilon)\ge2R+1.
\]
\end{corollary}

\begin{proof}
On the minus branch $p>m$.  Indeed, for $b=v_2(p+1)=39$ or $40$, the
coefficient-one candidates $2^{39}-1$ and $2^{40}-1$ are composite, so
primality forces the next odd coefficient and gives $p>2^{40}\ge m$; for
$b\ge41$ the inequality is immediate.  The odd-moment Newton exclusion in
characteristic greater than the support weight applies to (3.1), ruling out
every reduced signed polynomial of weight at most $2R$.
\end{proof}

\begin{remark}
The code in Theorem~\ref{thm:coupled} is one $\F_p$ subfield-subcode of an
extension-field GRS parity check.  Its $2^{40}$ singleton proportionality
classes are Frobenius-coupled.  Treating them as independent one-coordinate
codes would replace (3.2) by a false direct sum and would incorrectly
factorize its mass.
\end{remark}

\section{Weighted collision and Fourier identities}

\begin{theorem}[Pointwise weighted collision floor]\label{thm:floor}
Let $A:\F_p^m\to\F_p^d$ have rank $d$.  For
$S\subseteq\{1,\ldots,m\}$ put $\Phi(S)=A1_S$ and
\[
 N(v)=|\{S:\Phi(S)=v\}|.
\]
Then
\[
 2^m Z(A)=\sum_vN(v)^2,                                \tag{4.1}
\]
and, pointwise for every such map,
\[
 Z(A)\ge\max\{1,2^m/p^d\}.                            \tag{4.2}
\]
If $a_s$ denotes the $s$-th syndrome column, then
\begin{align}
 Z(A)
 &=p^{-d}\sum_{u\in\F_p^d}
   \prod_{s=1}^m
   \left(1+\cos\frac{2\pi\langle u,a_s\rangle}{p}\right) \tag{4.3}\\
 &=\frac{2^m}{p^d}\sum_{u\in\F_p^d}
   \prod_{s=1}^m
   \cos^2\frac{\pi\langle u,a_s\rangle}{p}.          \tag{4.4}
\end{align}
\end{theorem}

\begin{proof}
For an ordered collision $(S,T)$, remove the common intersection.  The
remaining signed incidence vector is a ternary kernel word $\epsilon$, and
the common intersection is arbitrary on its $m-\wt(\epsilon)$ zero
coordinates.  Summing gives (4.1).  Diagonal collisions give $Z\ge1$.
Cauchy--Schwarz over the $p^d$ syndromes gives
$\sum_vN(v)^2\ge4^m/p^d$, proving (4.2).

For (4.3), insert additive-character orthogonality for $A\epsilon=0$.
At one coordinate the weighted sum is
\[
 1+\tfrac12e^{2\pi i\theta/p}+\tfrac12e^{-2\pi i\theta/p}
 =1+\cos(2\pi\theta/p).
\]
The identity $1+\cos(2x)=2\cos^2x$ gives (4.4).
\end{proof}

\begin{corollary}[Minus-window mass statement]
For the code in Theorem~\ref{thm:coupled},
\[
 Z_W=2^{-m}\sum_vN_W(v)^2
 \ge\max\{1,2^m/p^{hR}\}.
\]
\end{corollary}

\begin{remark}[Normalization fence]
The factor $1+2\cos$ is the character factor for the unweighted count of
ternary words.  It is not the factor for the Prize-route mass
$\sum2^{-\wt}$.  Any mass argument using $1+2\cos$ has changed the object by
an exponential weight and must be re-derived from (4.3) or (4.4).
\end{remark}

\section{Consequences and nonclaims}

The packet supplies a corrected structural terminal:
\begin{center}
\begin{tabular}{@{}lll@{}}
\toprule
branch & proved code model & remaining task\\
\midrule
$p\equiv1\pmod4$ & bounded direct sum of prime-field GRS kernels
  & weighted mass upper bound\\
$p\equiv3\pmod4$ & one Frobenius-coupled root code (3.2)
  & weighted mass upper bound\\
\bottomrule
\end{tabular}
\end{center}

The common upper-bound target is a capacity-edge weighted subset-syndrome
$L^2$ estimate.  Theorem~\ref{thm:floor} is a lower bound and does not close
it.  Corollary~\ref{cor:distance} controls low-weight words but not the
weighted number of words above the distance floor.

This note does not claim:
\begin{itemize}
\item an upper bound on either branch's weighted mass;
\item a quotient-prefix maximum-fiber theorem;
\item payment of an MCA, CA, list, or adjacent-row numerator;
\item resolution of a Proximity Prize endpoint;
\item that lower order layers, parity normalization, or a first-match ledger
      are discharged by the top-window code theorem.
\end{itemize}

The companion script
\begin{center}
\small
\texttt{experimental/}\\
\texttt{f2\_admissible\_branch\_correction\_verify.py}
\end{center}
checks the M61 certificate, the dyadic order law, Frobenius-root cardinalities,
the weighted collision/Fourier normalization, and explicit mutation failures.

\end{document}
